\documentclass{baustms}
\citesort
\DeclareMathOperator{\Reop}{Re}

\theoremstyle{cupthm}
\newtheorem{thm}{Theorem}[section]
\newtheorem{prop}[thm]{Proposition}
\newtheorem{cor}[thm]{Corollary}
\newtheorem{lemma}[thm]{Lemma}
\numberwithin{equation}{section}

\newcommand{\Ass}{\mathrm{Ass}}
\newcommand{\HH}{\mathbb{H}}
\newcommand{\OO}{\mathbb{O}}
\newcommand{\RR}{\mathbb{R}}
\newcommand{\Sed}{\mathbb{S}}
\newcommand{\Imm}{\operatorname{Im}}
\newcommand{\rk}{\operatorname{rk}}
\newcommand{\one}[1]{\mathbf{1}[#1]}

\begin{document}
\runningtitle{Birth off the basis}
\title{Birth off the basis: the anti-associator and associator kernels of every quaternionic pair of sedenions}
\cauthor
\author[1]{Demetrios C. Agourakis}
\address[1]{Pontif\'icia Universidade Cat\'olica de S\~ao Paulo, Sorocaba, Brazil, and Faculdade S\~ao Leopoldo Mandic, Campinas, Brazil. ORCID 0009-0001-8671-8878\email{demetrios@agourakis.med.br}}
\authorheadline{D.~C. Agourakis}

\begin{abstract}
The anti-associator kernel of a pair of basis units in a Cayley--Dickson algebra was computed in earlier work. This article leaves the basis. A quaternionic pair in the sedenions lies in an octonion subalgebra over a quaternion subalgebra, and both kernels split along that subalgebra and its orthogonal complement. Each meets the subalgebra in dimension~$4$. Two explicit matrices over the quaternions govern the complement, and all six combinations occur: anti-associator dimension $4$ or $12$, associator dimension $4$, $8$ or $12$. For alternative elements the splitting iterates up the tower, each block a multiple of~$4$. One level higher the congruence modulo~$8$ fails. An explicit pair in dimension~$32$ has anti-associator kernel of dimension~$8$, and at that level the kernel has dimension $4$, $8$, $12$, $20$ or $28$.
\end{abstract}

\classification[2020]{primary 17A75; secondary 17D05}
\keywords{Cayley--Dickson algebra, sedenion, associator, quaternionic pair, alternative element}

\maketitle

\section{Setting and statement}

Write $\Sed=\OO\times\OO$ with the Cayley--Dickson product
\begin{equation}
(x,y)(a,b)=(xa-b\,\overline{y},\;\overline{x}\,b+ay),\qquad \overline{(x,y)}=(\overline{x},-y).
\label{eq:mul}
\end{equation}
An element $(a,b)$ is imaginary if and only if $a\in\Imm\OO$; $b$ is unrestricted. Write $b=b_0+\vec{b}$ for the real and imaginary parts of $b\in\OO$, and $u\times v$, $\langle u,v\rangle$ for the cross product and inner product of imaginary elements of a quaternion subalgebra, so that $uv=-\langle u,v\rangle+u\times v$ there. The associator is $[x,z,y]=(xz)y-x(zy)$, the anti-associator $W_{x,y}(z)=(xz)y+x(zy)$; $\Ass[x,y]=\ker[x,\cdot,y]$ and $\Ass'[x,y]=\ker W_{x,y}$, as in Biss--Dugger--Isaksen \cite{bdi} (BDI), Definition~8.1. A \emph{quaternionic pair} (BDI, Definition~4.6) is a pair of nonzero orthogonal imaginary elements with $[x,x,y]=0=[y,y,x]$.

\begin{thm}[birth at the sedenions, off the basis]\label{thm:1}
Let $(x,y)$ be a quaternionic pair in $\Sed$.
\begin{enumerate}
\item[\rm(i)] \emph{Containment.} There is a quaternion subalgebra $Q\subset\OO$ such that $x=(a,b)$, $y=(c,d)$ with $a,c\in\Imm Q$ and $b,d\in Q$. Orthogonality reads $\langle a,c\rangle+\langle b,d\rangle=0$.
\item[\rm(ii)] \emph{Anti-associator.} With $r:=[c,b]+[a,d]=2(c\times\vec{b}+a\times\vec{d})\in\Imm Q$,
\begin{equation}
\dim\Ass'[x,y]=4+8\cdot\one{r=0}.
\label{eq:N}
\end{equation}
\item[\rm(iii)] \emph{Associator.} Put $p=a\times c+\vec{b}\times\vec{d}$, $q=d_0\vec{b}-b_0\vec{d}$, $w=d_0 a-b_0 c$ (all in $\Imm Q$) and $t=\langle a,\vec{d}\rangle-\langle\vec{b},c\rangle\in\RR$, and let
\begin{equation}
M=\begin{pmatrix}p+q&w+t\\ w-t&p-q\end{pmatrix}
\label{eq:M}
\end{equation}
act on $Q\times Q$ by left multiplication, $(u,v)\mapsto\bigl((p+q)u+(w+t)v,\;(w-t)u+(p-q)v\bigr)$. Then
\begin{equation}
\dim\Ass[x,y]=4+4\bigl(2-\rk_{\HH} M\bigr)=4+\dim_{\RR}\ker M\in\{4,8,12\}.
\label{eq:A}
\end{equation}
In particular $\dim\Ass[x,y]=12$ if and only if $p=q=w=0$ and $t=0$; and $\rk_{\HH} M=1$ if and only if $M\neq 0$ and
$|p+q|^2|p-q|^2+|w+t|^2|w-t|^2=2\Reop\bigl(\overline{(p-q)}(w-t)\overline{(p+q)}(w+t)\bigr)$
(the Study determinant).
\end{enumerate}
\end{thm}

\begin{cor}[values and defect]\label{cor:1}
For every quaternionic pair of $\Sed$: $\dim\Ass'\in\{4,12\}$, $\dim\Ass\in\{4,8,12\}$, and all six combinations occur. The defect $\delta(x,y):=\dim\Ass[x,y]+\dim\Ass'[x,y]-16$ takes exactly the values $-8,-4,0,4,8$. A generic pair (Zariski-open in the pairs of $\Imm E$) has $(\dim\Ass',\dim\Ass)=(4,4)$, $\delta=-8$. The basis pairs, and every pair reached by BDI's lifting lemmas 8.7--8.10 from the octonions, have $\delta=0$.
\end{cor}

\begin{cor}[the mod-8 law off the basis]\label{cor:2}
$\dim\Ass'[x,y]\equiv 4\pmod{8}$ for every quaternionic pair of $\Sed$, not only for pairs of basis units. The companion statement for the associator fails: $\dim\Ass[x,y]=8$ occurs (for instance $x=e_1-e_8$, $y=e_2-e_9$).
\end{cor}

\begin{cor}[alternativity reading]\label{cor:3}
For $z=(a,b)\in\Imm\Sed$ let $\nu(z)=[a,b]\in\Imm\OO$. By Moreno \cite{moreno} (Theorem~3.3 with Proposition~3.1), $z$ is alternative if and only if $\nu(z)=0$. Since $\nu$ is quadratic, $r=\nu(x+y)-\nu(x)-\nu(y)$, so for a quaternionic pair of alternative elements, the setting of BDI Section~11,
\[
\dim\Ass'[x,y]=\begin{cases}12& \text{if $x+y$ is alternative,}\\ 4& \text{otherwise,}\end{cases}
\]
and by BDI Proposition~11.1 the annihilator of $x+y\,\ell'$ in $A_5$ has that dimension.
\end{cor}

\begin{cor}[propagation]\label{cor:4}
By BDI Lemmas 8.7--8.10, the four lifts of $(x,y)$ to $A_5$ have $(\dim\Ass',\dim\Ass)$ equal to $(2N+4,\,2A-4)$ for the same-side lifts and $(2A-4,\,2N+4)$ for the cross lifts, where $(N,A)$ are the values of Theorem~\ref{thm:1}; hence $\delta$ doubles at each lift and the pairs of $A_m$ obtained by lifting from $\Sed$ have defects $2^{m-4}\cdot\{-8,-4,0,4,8\}$.
\end{cor}

\section{Proof}

\begin{lemma}[containment]\label{lem:A}
If $x,y\in\Imm\Sed$ satisfy $[x,x,y]=0=[y,y,x]$, then $x,y\in Q\times Q$ for some quaternion subalgebra $Q\subset\OO$.
\end{lemma}

\begin{proof}
Let $x=(a,b)$, $y=(u,v)$ with $a\in\Imm\OO$. By \eqref{eq:mul}, $x^2=(a^2-b\overline{b},\,\overline{a}b+ab)=-(|a|^2+|b|^2)$, real. Also $xy=(au-v\overline{b},\,-av+ub)$ and, using alternativity of $\OO$ in the forms $a(au)=a^2 u$, $(ub)\overline{b}=u(b\overline{b})$, $(v\overline{b})b=v(\overline{b}b)$,
\[
x(xy)=\bigl(-(|a|^2+|b|^2)u+[a,v,\overline{b}],\; -(|a|^2+|b|^2)v+[a,u,b]\bigr).
\]
Hence $[x,x,y]=x^2 y-x(xy)=0$ if and only if $[a,u,b]=0$ and $[a,v,\overline{b}]=-[a,v,b]=0$: both components of $y$ lie in $\Ass_{\OO}[a,b]=\{z\in\OO:[a,z,b]=0\}$. When $\Imm a$ and $\Imm b$ are linearly independent this kernel is the quaternion subalgebra $Q(a,b)=\operatorname{span}\{1,a,b,ab\}$ (the associator depends only on the imaginary parts; after an orthogonalisation it is BDI Lemma~8.5, $\Ass_{\OO}=\HH$). So: if $x$ is \emph{nondegenerate} (imaginary parts of $a$, $b$ independent), $u,v\in Q(a,b)$ and $y\in Q(a,b)\times Q(a,b)$; if $y$ is nondegenerate too, symmetrically $a,b\in Q(u,v)$, so $Q(a,b)=Q(u,v)$. If $y$ is degenerate (imaginary parts of $u$, $v$ dependent), it lies in $Q(a,b)\times Q(a,b)$ anyway. If both are degenerate, the imaginary parts of $a,b,u,v$ span at most two dimensions and lie in a common $Q$. In every case $x,y\in Q\times Q$.
\end{proof}

\begin{lemma}[splitting]\label{lem:B}
Let $Q\subset\OO$ be a quaternion subalgebra, $E=Q\times Q$ and $E^\perp=Q^\perp\times Q^\perp$. For $x,y\in E$ the operators $W_{x,y}$ and $[x,\cdot,y]$ preserve $E$ and $E^\perp$; $E$ is the Cayley--Dickson double of $Q$, hence an octonion algebra; and for a quaternionic pair in $E$ both kernels meet $E$ in dimension exactly~$4$.
\end{lemma}

\begin{proof}
$Q\cdot Q^\perp\subset Q^\perp$ and $Q^\perp\cdot Q\subset Q^\perp$ in $\OO$, so \eqref{eq:mul} gives $E\cdot E^\perp\subset E^\perp$, $E^\perp\cdot E\subset E^\perp$, and $E\cdot E\subset E$. The product \eqref{eq:mul} restricted to $Q\times Q$ is the Cayley--Dickson double of $Q\cong\HH$, so $E\cong\OO$. In an octonion algebra, for orthogonal imaginary $x,y$, $\Ass[x,y]=\HH(x,y)$ and $\Ass'[x,y]=\HH(x,y)^\perp$ (BDI Lemma~8.5), each of dimension~$4$.
\end{proof}

\begin{lemma}[the $E^\perp$ blocks]\label{lem:C}
Choose $e\in Q^\perp$ a unit and write $Q^\perp=Qe$, so that $\OO=Q\times Q$ with \eqref{eq:mul} and $E^\perp=\{((0,u'),(0,v')):u',v'\in Q\}$. For $x=((a,0),(b,0))$, $y=((c,0),(d,0))$ with $a,c\in\Imm Q$, $b,d\in Q$, $\langle a,c\rangle+\langle b,d\rangle=0$,
\begin{align*}
W_{x,y}(u',v')&=(rv',\,-ru'),\\
[x,(u',v'),y]&=2\bigl((p+q)u'+(w+t)v',\;(w-t)u'+(p-q)v'\bigr)=2M(u',v'),
\end{align*}
with $r,p,q,w,t$ as in Theorem~\ref{thm:1}.
\end{lemma}

\begin{proof}
This is a direct expansion of \eqref{eq:mul} at the two levels. Writing $z=(U,V)$ with $U=(0,u')$, $V=(0,v')$ one finds $xz=((0,-au'-\overline{b}v'),(0,av'+bu'))$ and $zy=((0,cu'+\overline{d}v'),(0,-du'-cv'))$; multiplying out $(xz)y$ and $x(zy)$ and collecting, the coefficient of $u'$ in the first component of $W$ is $-(ac+ca)+(\overline{d}b+\overline{b}d)=2\langle a,c\rangle+2\langle b,d\rangle=0$, the coefficient of $v'$ in the second component is likewise $0$, and the two remaining coefficients are $[\overline{b},c]+[\overline{d},a]=[c,b]+[a,d]=r$ and its negative. For the associator the four coefficients are $[a,c]+2\Imm(\overline{d}b)$, $\{\overline{d},a\}-\{\overline{b},c\}$, $\{a,d\}-\{b,c\}$ and $[a,c]+2\Imm(d\overline{b})$; expanding $b=b_0+\vec{b}$, $d=d_0+\vec{d}$ gives $2(p+q)$, $2(w+t)$, $2(w-t)$, $2(p-q)$.
\end{proof}

\begin{proof}[Proof of Theorem~\ref{thm:1}]
(i) is Lemma~\ref{lem:A}. By Lemma~\ref{lem:B}, $\ker W_{x,y}=(\ker W_{x,y}\cap E)\oplus(\ker W_{x,y}\cap E^\perp)$, the first summand of dimension~$4$, and likewise for the associator. By Lemma~\ref{lem:C} the restriction to $E^\perp$ is left multiplication on $Q\times Q\cong\HH^2$ by a $2\times 2$ quaternion matrix; left multiplication commutes with right scalar multiplication by $Q$, so each kernel is a right $Q$-subspace of $Q^2$ and has real dimension $0$, $4$, or $8$. For $W$ the matrix is $\bigl(\begin{smallmatrix}0&r\\-r&0\end{smallmatrix}\bigr)$: kernel $Q^2$ if $r=0$, and $0$ if $r\neq 0$. This is \eqref{eq:N}. For the associator the matrix is $M$ and \eqref{eq:A} is the statement that $\dim_{\RR}\ker M=4(2-\rk_{\HH} M)$; the Study-determinant form of the rank-one condition is the usual one for $2\times 2$ quaternion matrices, $|\alpha|^2|\delta|^2+|\beta|^2|\gamma|^2-2\Reop(\overline{\delta}\,\gamma\,\overline{\alpha}\,\beta)=|\alpha|^2|\delta-\gamma\alpha^{-1}\beta|^2$ when $\alpha\neq 0$.
\end{proof}

\begin{proof}[Proof of the corollaries]
The value sets are read off \eqref{eq:N} and \eqref{eq:A}. All six combinations occur: $(4,4)$ at $(e_1-e_2,e_3-e_9)$; $(4,8)$ at $x=-e_2+e_3-e_9+e_{10}$, $y=-e_1-e_3-e_9-e_{11}$; $(4,12)$ at $(e_1,e_{10})$; $(12,4)$ at $(e_1,e_2)$; $(12,8)$ at $(e_1-e_8,e_2-e_9)$; $(12,12)$ at $(e_1-e_{10},e_1+e_{10})$. Genericity: $r\neq 0$ and $\rk_{\HH} M=2$ are open conditions and hold at the first witness. The doubling itself is immediate: $(2N+4)+(2A-4)-32=2(N+A-16)$. That BDI's Lemmas 8.7--8.10 send $(N,A)$ to $(2N+4,2A-4)$ or $(2A-4,2N+4)$ is their statement, checked in this convention (item~(7) below); together with $\delta=0$ at Lemma~8.5 this gives Corollary~\ref{cor:4} and the last sentence of Corollary~\ref{cor:1}. Corollary~\ref{cor:2} is immediate. For Corollary~\ref{cor:3}, Moreno's Theorem~3.3 says that a pure $(a,b)$ with $a$, $b$ pure is alternative if and only if $a$, $b$ are alternative (automatic in $\OO$) and linearly dependent, and his Proposition~3.1 removes the real part of $b$; together, $(a,b)$ with $a$ pure is alternative if and only if $[a,b]=0$. Bilinearity gives $r=\nu(x+y)-\nu(x)-\nu(y)$.
\end{proof}

\section{Remarks}

\paragraph{Where the basis sits.}
For basis units $b,d\in\{0\}\cup\{\pm e_k\}$ and $a,c\in\{0\}\cup\{\pm e_j\}$ with at most one of $a$, $b$ and one of $c$, $d$ nonzero. One checks that then $\rk_{\HH} M\in\{0,2\}$ and that $\dim\Ass=16-\dim\Ass'$: the complement identity $\Ass\oplus\Ass'=\Sed$ used to write the differing clause as $2^{n+1}-2N-4$ is the statement $\delta=0$, and Theorem~\ref{thm:1} shows it is a property of the defect-zero stratum, not of quaternionic pairs in general.

\paragraph{What decides $\Ass'$.}
Within a fixed plane $\Pi=\operatorname{span}\{x,y\}$ the associator kernel is constant (the associator is bilinear and antisymmetric in $x,y$) but $\Ass'$ is not: in $\HH(e_1,e_{10})$ the pair $(e_1,e_{10})$ has $\dim\Ass'=4$ while the pair $(e_1-e_{10},e_1+e_{10})$ in the same plane has $12$. Theorem~\ref{thm:1} locates the jump: $r$ is a symmetric bilinear form on $\Imm E$ with values in $\Imm Q$, $r(z,z)=2\nu(z)$, and $\Ass'$ is large exactly on its isotropic pairs. Moreno's classification of alternative elements, which does not mention kernels, is therefore the kernel theory of $W$ at the sedenions.

\paragraph{Positive defect.}
If $\delta>0$ then $\Ass\cap\Ass'\neq 0$, and any nonzero $z$ in the intersection satisfies $(xz)y=0=x(zy)$; so $xz$ is a zero divisor or zero. Negative defect has no such reading; it is the generic case.

\paragraph{What this does not do.}
It does not give the birth function for quaternionic pairs of $A_5$ and beyond that are not lifts: Lemma~\ref{lem:A} uses alternativity of the components, which fails in $\Sed$. Section~\ref{sec:tower} recovers the containment, and a block decomposition, for pairs of alternative elements at every level; pairs with a non-alternative member at $m\ge 5$ remain not established. It does not describe the Zariski strata of $\rk_{\HH} M=1$ in invariant terms beyond the Study determinant. Nothing here is formalized: the development in Lean of the basis computation does not contain alternativity of $A_3$, which Lemmas~\ref{lem:A} and~\ref{lem:B} need.

\section{Verification}

The tools named in the acknowledgement were used to propose candidate errors in this argument. That reading is not a step of the proof. What follows is exact integer arithmetic (the multiplication \eqref{eq:mul} by structural recursion, ranks modulo the primes $1\,000\,003$ and $1\,048\,573$).
(1)~All $7\,413$ quaternionic pairs of weight two: Lemma~\ref{lem:A} holds for every pair; \eqref{eq:N} and \eqref{eq:A} agree with the direct $16\times 16$ nullities for every pair; the values are $(4,4)\times 4200$, $(4,12)\times 924$, $(12,4)\times 1827$, $(12,8)\times 378$, $(12,12)\times 84$.
(2)~$20\,000$ random integer pairs in $E$ for $Q=\HH(e_1,e_2)$ with entries in $[-3,3]$: \eqref{eq:N} and \eqref{eq:A} agree in every case; values $(4,4)\times 19679$, $(4,8)\times 56$, $(4,12)\times 5$, $(12,4)\times 233$, $(12,8)\times 26$, $(12,12)\times 1$. The combination $(4,8)$, absent at weight two, appears.
(3)~$60\,000$ further pairs: $\dim\ker M=8$ exactly when $p=q=w=t=0$.
(4)~The octonion fact of Lemma~\ref{lem:A}, $\dim\Ass_{\OO}[a,b]=4$ for independent imaginary parts: $3000/3000$ random pairs.
(5)~Moreno's criterion in the convention \eqref{eq:mul}: alternative if and only if $[a,b]=0$, $4000/4000$ random imaginary elements, $2007$ of them alternative.
(6)~Inside the coordinate quaternion $\HH(e_1,e_{10})$: of the $546$ primitive integer orthogonal pairs with entries in $[-3,3]$, $500$ have $(4,4)$, $36$ have $(4,12)$, $8$ have $(12,4)$, $2$ have $(12,12)$, exactly the strata $r=0$ and $\rk_{\HH} M$ predict.
(7)~Corollary~\ref{cor:4}'s lift rule: the six witness pairs and $400$ random quaternionic pairs of $E$, all four lifts each, $406/406$; every lift is again a quaternionic pair.
(8)~Lemma~\ref{lem:D}: alternative if and only if (I) and (II), on $150$ elements of $A_5$ and $150$ of $A_6$ drawn from three families, $300/300$; property (N), dimension $2$, for $1000$ dense and $1008$ sparse pure elements of $A_4$, $300+329$ of $A_5$, $60+123$ of $A_6$.
(9)~The identity $[u,b,b]=-[b,b,u]$ cited in Lemma~\ref{lem:D}, and flexibility $[b,u,b]=0$: $400/400$ random elements in each of $A_3$, $A_4$, $A_5$.

\noindent\textit{Status of Section~\ref{sec:tower} is stated there.}
Theorem~\ref{thm:1} and Corollaries~\ref{cor:1}--\ref{cor:3} are proved above and checked on every step. For Corollary~\ref{cor:4} the doubling identity is proved here; the lift map is BDI's, checked in this convention on $406$ pairs. The literature read for this article is BDI \cite{bdi} and Moreno \cite{moreno}. Later work on sedenion zero divisors was not searched, and no priority claim is made.

\section{The tower, for pairs of alternative elements}\label{sec:tower}

Lemma~\ref{lem:A} does not survive the sedenions as stated, because its proof uses alternativity of the components. For the class of pairs that BDI's Section~11 is about, quaternionic pairs of alternative elements, Moreno's classification supplies the containment at every level, and with it a decomposition of both kernels along a tower of doubled subalgebras.

\begin{lemma}[alternative elements in the convention \eqref{eq:mul}]\label{lem:D}
Let $n\ge 3$ and $x=(a,b)\in A_{n+1}$ with $a\in\Imm A_n$, $b\in A_n$. Then $x$ is alternative if and only if
\begin{equation}
\text{\rm(I)}\quad [a,u,b]=0\qquad\text{and}\qquad\text{\rm(II)}\quad [a,a,u]+[b,b,u]=0
\label{eq:I}
\end{equation}
for all $u\in A_n$. Suppose moreover that $A_n$ satisfies {\rm(N)}: for every nonzero pure $a$, $\{b:[a,u,b]=0\text{ for all }u\}=\RR\oplus\RR a$. Then $x$ is alternative if and only if either $a\neq 0$ is alternative and $b\in\RR\oplus\RR a$, or $a=0$ and $\Imm b$ is alternative. Property {\rm(N)} holds at $n=3$ (BDI Lemma~8.5), at $n=4$ by Theorem~\ref{thm:1}, and for every $n\ge 4$ it is Moreno's Theorem~2.3 \cite{moreno}.
\end{lemma}

\begin{proof}
$x^2=-(|a|^2+|b|^2)$ is real. For $z=(u,v)$ the raw expansion in the proof of Lemma~\ref{lem:A}, before alternativity of $\OO$ was used, is
\[
x(xz)=\bigl(a(au)-(ub)\overline{b}+[a,v,\overline{b}],\; a(av)-(v\overline{b})b+[a,u,b]\bigr).
\]
So $[x,x,z]=0$ for all $z$ if and only if, separating $u$ and $v$, $[a,v,\overline{b}]=0$ and $[a,u,b]=0$ for all $u,v$, and $a(au)-(ub)\overline{b}=-|x|^2 u$, $a(av)-(v\overline{b})b=-|x|^2 v$. Since $\overline{b}=2b_0-b$ and the associator kills real entries, the first-order conditions are (I). With $a^2=-|a|^2$, $b\overline{b}=|b|^2$ and flexibility of Cayley--Dickson algebras ($[u,b,b]=-[b,b,u]$; Moreno, Section~1), $a(au)+|a|^2 u=-[a,a,u]$ and $(ub)\overline{b}-|b|^2 u=[u,b,\overline{b}]=-[u,b,b]=[b,b,u]$, so the second-order conditions are (II). Given (N): if $a\neq 0$, (I) gives $b=s+ta$ and (II) becomes $(1+t^2)[a,a,u]=0$; if $a=0$, (II) is $[b,b,u]=0$. Conversely both families satisfy (I) ($[a,u,a]=0$ by flexibility) and (II). For (N) at $n=4$: let $a\neq 0$ be pure in $\Sed$ and $[a,u,b]=0$ for all $u$; put $b'=\Imm b-(\langle\Imm b,a\rangle/|a|^2)a$, so $[a,u,b']=0$ for all $u$. If $b'\neq 0$ then $(a,b')$ is a pair of nonzero orthogonal imaginary elements with $[a,a,b']=0$ ($u=a$) and $[b',b',a]=-[a,b',b']=0$ ($u=b'$): a quaternionic pair with $\Ass[a,b']=\Sed$, of dimension~$16$, contradicting the bound~$12$ of Theorem~\ref{thm:1}(iii). Hence $b'=0$.
\end{proof}

\begin{thm}[tower decomposition]\label{thm:2}
Let $m\ge 3$ and let $(x,y)$ be a quaternionic pair of alternative elements of $A_m$. Then:
\begin{enumerate}
\item[\rm(i)] there is a quaternion subalgebra $Q\subset\OO=A_3$ such that $x,y\in E_m:=Q^{2^{m-3}}$, the $(m-3)$-fold Cayley--Dickson double of $Q$ inside $A_m$; $E_m$ is a subalgebra isomorphic to $A_{m-1}$;
\item[\rm(ii)] $E_m^\perp=(Q^\perp)^{2^{m-3}}$ is invariant under $W_{x,y}$ and $[x,\cdot,y]$, so
\begin{align*}
\dim\Ass'_{A_m}[x,y]&=\dim\Ass'_{E_m}[x,y]+\beta_m,\\
\dim\Ass_{A_m}[x,y]&=\dim\Ass_{E_m}[x,y]+\alpha_m,
\end{align*}
where $\beta_m,\alpha_m$ are the kernel dimensions of the restrictions to $E_m^\perp$; both restrictions commute with the right action of $Q$ on $(Q^\perp)^{2^{m-3}}\cong\HH^{2^{m-3}}$, so $\alpha_m,\beta_m\in 4\mathbb{Z}$ and $0\le\alpha_m,\beta_m\le 2^{m-1}$;
\item[\rm(iii)] under an isomorphism $E_m\cong A_{m-1}$ the pair stays a quaternionic pair of alternative elements, so {\rm(i)--(ii)} iterate down to $E_4\cong A_3=\OO$, where $\Ass'=\HH(x,y)^\perp$ and $\Ass=\HH(x,y)$ (BDI Lemma~8.5): with $\beta_k,\alpha_k$ the kernels of the restrictions to the complement of $E_k$ in $A_k$, transported along the isomorphisms,
\[
\dim\Ass'[x,y]=4+\sum_{k=4}^{m}\beta_k,\qquad
\dim\Ass[x,y]=4+\sum_{k=4}^{m}\alpha_k.
\]
At $k=4$, $\beta_4=8\cdot\one{r=0}$ and $\alpha_4=4(2-\rk_{\HH} M)$ by Theorem~\ref{thm:1}.
\end{enumerate}
\end{thm}

\begin{proof}
By Lemma~\ref{lem:D}, equivalently Moreno \cite{moreno}, Theorem~3.3 with Proposition~3.1, for $n\ge 3$ a pure element of $A_{n+1}$ is alternative if and only if it has the form $(p,s+tp)$ with $p$ a pure alternative element of $A_n$ and $s,t\in\RR$, or the form $(0,s+p)$. For $n=3,4$ this is proved above from BDI Lemma~8.5 and Theorem~\ref{thm:1}; for $n\ge 5$ it imports property (N) from Moreno's Theorem~2.3. Hence an alternative $x\in A_m$ lies in $C_p\times C_p$ with $C_p=\operatorname{span}\{1,p\}$, and $y$ in $C_{p'}\times C_{p'}$; both lie in $B\times B$ for any subalgebra $B\subset A_{m-1}$ containing $p,p'$. Now $p,p'$ are themselves pure alternative elements of $A_{m-1}$, so they lie in doubles of the same shape over pure alternative elements of $A_{m-2}$, and $B$ may be taken of the form $B_1\times B_1$. Descending to $A_3=\OO$, where every element is alternative, we reach pure imaginary octonions and may take any quaternion subalgebra containing them. Unwinding, $x,y\in Q^{2^{m-3}}=E_m$. This is (i). For (ii): $Q\cdot Q^\perp\subset Q^\perp$ and $Q^\perp\cdot Q\subset Q^\perp$ in $\OO$, and \eqref{eq:mul} only ever multiplies coordinates, so $E_m\cdot E_m^\perp\subset E_m^\perp$ and $E_m^\perp\cdot E_m\subset E_m^\perp$; the restriction of $W_{x,y}$ to $E_m^\perp$ is built from left and right multiplications by elements of $E_m$, which act on each $Q^\perp$-coordinate by left multiplications by elements of $Q$, and these commute with right multiplication by $Q$. (iii) is immediate from (i), (ii), and Theorem~\ref{thm:1}.
\end{proof}

\begin{cor}[the bin digits are blocks]\label{cor:5}
For a pair of basis units of $A_m$, $m\ge 5$, define $\varepsilon_k\in\{0,1\}$ by writing the basis birth number as $4+\sum_{k=4}^{m}2^{k-1}\varepsilon_k$. At $m=5$, for all $465$ pairs, $\beta_k=2^{k-1}\varepsilon_k$ and $\alpha_k=2^{k-1}(1-\varepsilon_k)$, $k=4,5$: each block of the tower is either $0$ or the whole complement, and the digit at position $k$ records which. This is observed at $m=5$, exhaustively on the basis. The dynamics of those digits for general $m$ is not established here.
\end{cor}

On $4002$ random quaternionic pairs of alternative elements of $A_5$ (Moreno's parametrisation, integer parameters in $[-2,2]$), containment in $Q^4$, invariance of $E_5^\perp$, and additivity of both kernels hold in every case. The values seen there are $(\dim\Ass',\dim\Ass)=(4,4)$ for $3215$ pairs, $(12,4)$ for $481$, $(28,4)$ for $300$, $(4,12)$ for $4$, $(12,12)$ for $1$, $(4,28)$ for $1$, with $\beta_5\in\{0,8,16\}$. That sample missed a stratum. The generic pair at level~$5$ still has $(4,4)$, defect $-24$.

\begin{prop}[the mod-8 law stops at the sedenions]\label{prop:6}
Theorem~\ref{thm:2} gives $\beta_m\in 4\mathbb{Z}$. The stronger congruence $\beta_m\in 8\mathbb{Z}$, and with it $\dim\Ass'\equiv 4\pmod{8}$ for every quaternionic pair of alternative elements, is false at $m=5$. The elements
\[
x=e_3-e_{16}-2e_{19},\qquad y=e_2-e_3-2e_{10}+2e_{11}-e_{16}
\]
of $A_5$ are alternative, orthogonal, and anticommute, and
\[
\dim\Ass'[x,y]=8,\qquad \dim\Ass[x,y]=4,\qquad \beta_5=4.
\]
The $4$ in the top block is the kernel of a $2\times 2$ Hermitian matrix over $\HH$ of Study determinant zero and not itself zero; right $\HH$-linearity accounts for the factor $4$ and does not double it.
\end{prop}

\begin{proof}
Alternativity, orthogonality, and the anticommutator are finite checks in the product \eqref{eq:mul}. The three dimensions are the nullities of integer matrices, computed by Gaussian elimination modulo $1\,000\,003$ and modulo $1\,048\,573$, which agree. In the notation of Theorem~\ref{thm:3} below: $p=-4e_1$, $a=c=4e_1$, $b=8$, so
\[
P=\begin{pmatrix}0&-4e_1\\ 4e_1&0\end{pmatrix},\qquad
B=\begin{pmatrix}4e_1&8\\ -8&4e_1\end{pmatrix};
\]
$P$ is invertible with $P^2=16$, and the Schur complement $S=P+BP^{-1}B$ equals $-16\bigl(\begin{smallmatrix}1&-e_1\\ e_1&1\end{smallmatrix}\bigr)$, whose kernel is $\{(e_1 g,\, g):g\in\HH\}$, a right $\HH$-line, of real dimension~$4$.
\end{proof}

\begin{prop}[the top block at level $5$]\label{prop:7}
Let $(x,y)$ be a quaternionic pair of alternative elements of $A_5$. Then $\beta_5\in\{0,4,8,16\}$. All four values occur. In particular $12$ does not.
\end{prop}

\begin{proof}
By Theorem~\ref{thm:2} the pair lies in $E=Q^4$ for some quaternion subalgebra $Q\subset\OO$. An automorphism of $\OO$ carries $Q$ to $\HH(e_1,e_2)$ and extends to $A_5$ by acting on each octonion coordinate; it preserves $W$ and the top doubling, so it is enough to treat $E=\operatorname{span}\{e_i:i\in\{0,1,2,3,8,9,10,11,16,17,18,19,24,25,26,27\}\}$. On the complement $F\times F$, with $F\cong\HH^2$, the operator $W_{x,y}$ is the anticommutator of the two $4\times 4$ matrices over $\HH$ that the doubling formula assigns to left multiplication by $x$ and right multiplication by $y$ (the identity is bilinear, and it holds on all $256$ pairs of basis elements of $E$). For an orthogonal pair with no component along $1$, that anticommutator has the shape $\bigl(\begin{smallmatrix}P&B\\ -B&P\end{smallmatrix}\bigr)$ with $P=\bigl(\begin{smallmatrix}0&p\\ -p&0\end{smallmatrix}\bigr)$, the operator of Theorem~\ref{thm:3}: the shape holds on every pair of distinct imaginary basis elements, and on each square $e_i^2$ the failure is exactly $2I$, which cancels against $\langle x,y\rangle=0$. If $p\neq 0$ then $P$ is invertible and the kernel is the kernel of a single $2\times 2$ matrix over $\HH$, hence of real dimension $0$, $4$, or $8$. If $p=0$ the kernel is a square of the kernel of the off-diagonal block, hence of real dimension $0$, $8$, or $16$. The four values occur: $0$, $8$, and $16$ in the level-$5$ census, and $4$ by Proposition~\ref{prop:6}.
\end{proof}

\section{Level $5$ closed}

By Theorem~\ref{thm:2} and the $G_2$ reduction of Proposition~\ref{prop:7} we may take $Q=\HH(e_1,e_2)$, so the complement of $E_5$ in $A_5$ is $F\times F$ with $F=Q^\perp\times Q^\perp$. Identify each copy of $Q^\perp$ with $\HH$ by $e_{4+i}\leftrightarrow e_i$ ($i=0,1,2,3$). The operators below are $2\times 2$ matrices over $\HH$ acting on columns by left multiplication of the entries. By Moreno's parametrisation an alternative pure element of $A_5$ lying in $E_5$ is
\begin{equation}
x=(\alpha,\, s+t\alpha),\qquad \alpha=(X_0,\,\sigma+\tau X_0),\qquad X_0\in\Imm Q,\quad s,t,\sigma,\tau\in\RR,
\label{eq:nf}
\end{equation}
and likewise $y$ with $Y_0,s',t',\sigma',\tau'$. Write $[u,v]=uv-vu$ and $\langle u,v\rangle$ for the inner product of imaginary quaternions.

\begin{thm}[the top block at level $5$, closed form]\label{thm:3}
Let $(x,y)$ be a quaternionic pair of alternative elements of $A_5$ in the normal form \eqref{eq:nf}. Then:
\begin{enumerate}
\item[\rm(i)] $\dim\Ass'_{E_5}[x,y]\in\{4,12\}$, and it is $12$ exactly when $t=t'$ or $\alpha$ and $\alpha'$ are parallel.
\item[\rm(ii)] On $F\times F$, $W_{x,y}(f,g)=(Pf+Bg,\,-Bf+Pg)$ with
\begin{align*}
P&=\begin{pmatrix}0&p\\ -p&0\end{pmatrix},&
B&=\begin{pmatrix}a&b\\ -\overline{b}&c\end{pmatrix},\\
p&=(1+tt')(\tau'-\tau)[X_0,Y_0],\\
b&=(t'-t)\bigl(2(\tau'-\tau)\langle X_0,Y_0\rangle+2\sigma' X_0-2\sigma Y_0\bigr),\\
a&=(t'-t)\bigl((1+\tau\tau')[X_0,Y_0]+2(\sigma'\tau X_0-\sigma\tau' Y_0)\bigr),\\
c&=(t'-t)\bigl((1+\tau\tau')[X_0,Y_0]-2(\sigma'\tau X_0-\sigma\tau' Y_0)\bigr).
\end{align*}
\item[\rm(iii)] $\beta_5=\dim_{\RR}\ker(P+BP^{-1}B)\in\{0,4,8\}$ if $p\neq 0$, and $\beta_5=2\dim_{\RR}\ker B\in\{0,8,16\}$ if $p=0$.
\item[\rm(iv)] If $\dim\Ass'_{E_5}[x,y]=12$ then $B=0$ and $\beta_5=16\cdot\one{p=0}$; for parallel $\alpha,\alpha'$ this gives $\beta_5=16$.
\item[\rm(v)] Consequently $\dim\Ass'[x,y]\in\{4,8,12,20,28\}$. All five values occur; $16$ and $24$ do not.
\end{enumerate}
\end{thm}

\begin{proof}
(i) is Theorem~\ref{thm:1} inside $E_5\cong A_4$: the pair is $(\alpha,s+t\alpha)$, $(\alpha',s'+t'\alpha')$ with $\alpha,\alpha'$ pure in the octonion algebra $E_4=Q\times Q$, so $r=(t-t')[\alpha',\alpha]$, which vanishes if and only if $t=t'$ or $[\alpha,\alpha']=0$, and two pure octonions commute if and only if they are parallel.

(ii) The map $(x,y)\mapsto W_{x,y}|_{F\times F}$ is bilinear, and so are the right-hand sides once $x,y$ are written in the sixteen $Q$-coordinates $(X_0,X_1,X_2,X_3)$, $(Y_0,\ldots,Y_3)$ of $E_5=Q^4$:
\begin{align*}
p&=[X_0,Y_1]-[X_1,Y_0]+[X_2,Y_3]-[X_3,Y_2],\\
a,c&=K\pm J,\\
K&=[X_0,Y_2]+[X_1,Y_3]-[X_2,Y_0]-[X_3,Y_1],\\
J&=(X_1\overline{Y_3}-\overline{X_1}Y_3)-(X_3\overline{Y_1}-\overline{X_3}Y_1),\\
b&=[X_0,Y_3]-[X_1,Y_2]-[X_2,Y_1]+[X_3,Y_0]\\
&\quad+2X_0\overline{Y_3}-X_2\overline{Y_1}-\overline{X_1}Y_2+2\overline{X_3}Y_0+\overline{X_1}\,\overline{Y_2}+\overline{X_2}\,\overline{Y_1},
\end{align*}
and the lower-left entry of $B$ is $-\overline{b}$. A bilinear identity holds if and only if it holds on pairs of basis vectors; the block shape, the left-multiplication form of every block, and these four formulas were checked on all $15\times 15$ pairs of imaginary basis units of $E_5$ (on the diagonal pairs $e_i,e_i$ the shape fails by exactly $2I$ in $P$, which is cancelled by $\langle x,y\rangle=0$, as in Proposition~\ref{prop:7}). Substituting $X_1=\sigma+\tau X_0$, $X_2=s+t X_0$, $X_3=t X_1$ and the primed analogues gives the displayed forms. The same identities were rechecked on $6000$ further alternative pairs, as were (i), (iv), additivity, and the value of $\beta_5$ by direct nullity.

(iii) $P^2=|p|^2 I$ since $p$ is a commutator, hence pure. If $p\neq 0$, the kernel equations $Pf+Bg=0=-Bf+Pg$ give $f=-P^{-1}Bg$ and then $(P+BP^{-1}B)g=0$; the kernel is the graph of $-P^{-1}B$ over $\ker(P+BP^{-1}B)$, a right $\HH$-submodule of $\HH^2$. The order matters: $P+P^{-1}BPB$ has a different kernel unless $P$ and $B$ commute, and gives $0$ on the pair of Proposition~\ref{prop:6}. Direct comparison of $\dim\ker(P+BP^{-1}B)$ with the nullity of $W$ confirms the formula on that pair and on $14\,120$ random pairs with $p\neq 0$. If $p=0$ the equations decouple into $Bg=0=Bf$.

(iv) If $t=t'$, every entry of $B$ carries the factor $t'-t$. If $\alpha'=\lambda\alpha$ then $Y_0=\lambda X_0$, $\sigma'=\lambda\sigma$, $\tau'=\tau$, so $[X_0,Y_0]=0$ and $\sigma'\tau X_0-\sigma\tau' Y_0=0$, whence $p=a=b=c=0$. With $B=0$ the operator is $P\oplus P$, invertible if $p\neq 0$ and zero if $p=0$.

(v) From (i), (iii), (iv): if $\dim\Ass'_{E_5}=12$ the total is $12$ or $28$; if it is $4$ the total is $4+\beta_5\in\{4,8,12,20\}$. Realizations: $4$ is generic; $8$ is Proposition~\ref{prop:6}; $12$ occurs both as $12+0$ and as $4+8$; $20$ and $28$ occur on the basis ($84$ and $171$ pairs, Corollary~\ref{cor:5}).
\end{proof}

\paragraph{Remark.}
Formulas (ii) and statements (iii)--(iv) do not use alternativity: they hold for every orthogonal imaginary pair of $E_5=Q^4$ (checked on $6000$ such pairs with arbitrary $X_1,Y_1$). Among those, $(\dim\Ass'_{E_5},\beta_5)=(4,16)$ occurs, for instance
\begin{align*}
x&=-3e_1+2e_3-2e_9-3e_{11}+3e_{17}-2e_{19}+2e_{25}+3e_{27},\\
y&=-2e_1-2e_9-e_{11}+2e_{16}-2e_{17}-2e_{25}-e_{27},
\end{align*}
with $\dim\Ass'=20$, $\dim\Ass=12$, neither element alternative. What alternativity buys is the containment in Theorem~\ref{thm:2}, not the block algebra.

\paragraph{Observed at levels $6$ and $7$.}
Moreno's parametrisation iterated, integer parameters in $[-2,2]$, standard $Q$. Level~$6$, $2000$ pairs: invariance of $E_6^\perp$, additivity of both kernels, and the block shape $W(f,g)=(Pf+Bg,-Bf+Pg)$ on the top complement ($P,B$ now $4\times 4$ over $\HH$) hold in every case; $\beta_6\in\{0,8,16,32\}$ ($1801$, $8$, $33$, $158$), $\alpha_6\in\{0,8,16,32\}$; $(\dim\Ass'_{E_6},\beta_6)=(12,8)$ occurs, so $\dim\Ass'=20$ arises at level~$6$ with a quarter block ($8$ of $32$). Level~$7$, $400$ pairs: the same structural checks all pass; $\beta_7\in\{0,32,64\}$, $\alpha_7=0$ throughout. The associator does not have the $W$ shape on the top complement once $s$ or $s'$ is nonzero, so no analogue of Theorem~\ref{thm:3} for $\alpha_k$ is offered here.

\paragraph{Status.}
Theorem~\ref{thm:2}(i)--(iii) is established. At $m=5$ it is self-contained (Lemma~\ref{lem:D} with property (N) from BDI Lemma~8.5 and Theorem~\ref{thm:1}); for $m\ge 6$ the single import is Moreno's Theorem~2.3 at $n\ge 5$, checked in this convention on $812$ elements of $A_5$ and $A_6$. Corollary~\ref{cor:5} is observed at $m=5$, on the basis. Propositions~\ref{prop:6} and~\ref{prop:7} and Theorem~\ref{thm:3} are established, so $\beta_5$ is explicit and the value set of $\dim\Ass'$ at level~$5$ is closed. At levels $6$ and $7$ the block shape and $\beta_k\in 4\mathbb{Z}$ are theorems; the value sets are observed only. The explicit form of $\beta_k$ for $k\ge 6$ and of $\alpha_k$ for $k\ge 5$ is not established. Quaternionic pairs with a non-alternative member at levels $\ge 5$ are not treated beyond the Remark.

\ack
Generative artificial intelligence tools were used on 3~October 2026 to draft and revise the prose of this article and to propose possible errors in the algebra. The tools were Grok~4.7 (xAI), used through the Cursor editor (\url{https://x.ai}), and, through a private OpenAI-compatible router in the author's laboratory, GLM-5.3 (Z.AI), Qwen3.8-27B and Goedel-Prover-V2-32B. The laboratory router is not a public service. Every algebraic objection was checked by the author against the product \eqref{eq:mul} by exact integer arithmetic, and none survived. The statements, the proofs and the decision of what is established are the author's. No generative model is an author. The C++ programs that perform the exact checks in the Verification section are available from the author and can be supplied as supplementary material if the editor wishes. There was no funding.

\begin{thebibliography}{9}
\bibitem{bdi}
D.~K. Biss, D.~Dugger, and D.~C. Isaksen,
Large annihilators in Cayley--Dickson algebras,
\emph{Communications in Algebra} \textbf{36} (2008), no.~2, 632--664.

\bibitem{moreno}
G.~Moreno,
Alternative elements in the Cayley--Dickson algebras,
arXiv:math/0404395, 2004.
\end{thebibliography}

\end{document}
